\section{Background \& State of the Art}

\subsection{Bacteriophage Host Prediction Problem}
\noindent - What bacteriophages are: viruses that infect bacteria, most abundant biological entities on the planet\\
- The host range property: each phage encodes structural proteins that determine which bacterial strains it can infect \\
- Why host specificity is both an advantage and a limitation for phage therapy\\
- The computational problem: given a phage genome and a set of candidate hosts $\mathcal{H} = \{h_1, \ldots, h_m\}$, identify the subset $\subseteq \mathcal{H}$ of true hosts\\
- Scale: hundreds of thousands of phage genomes, thousands of candidate hosts, experimental screening infeasible



% The problem of predicting which bacterial hosts a phage can infect from genomic data has been studied computationally 
% for over a decade. Early approaches relied on sequence alignment between phage genomes and bacterial genomes, exploiting 
% the co-evolutionary signal between phage tail proteins and bacterial surface receptors. These methods performed 
% reasonably well when closely related phage-host pairs existed in reference databases, but struggled with novel phages 
% that had no known close relatives.

% More recent approaches framed the problem as a supervised classification task. Given a set of experimentally verified 
% phage-host pairs, a classifier is trained to predict whether a new phage infects a given host based on features derived
% from the phage genome. Among the most informative features are functional annotations of phage proteins: the presence 
% and predicted function of proteins such as lysins, tail fibres, capsid proteins, and portal complexes, since these 
% proteins directly mediate host recognition and cell lysis.

% In this work, the predictions are produced at the protein level. For each protein in a phage genome and each candidate host
% h∈Hh∈H, a classifier outputs a score s∈[0,1]s∈[0,1] reflecting the probability that the protein's function is consistent 
% with infection of that host. Since a phage genome typically contains many proteins — ranging from roughly 10 to 50 in our
% dataset — and each protein may be annotated with one or more functional categories, this yields a structured collection of 
% scores indexed by (protein, host, function) triples. These protein-level scores are then aggregated to the phage level via 
% NaN-aware mean pooling: for each (phage, host, function) combination, we take the mean score across all proteins of that 
% phage annotated with that function, ignoring proteins with missing annotations for that function rather than treating them 
% as zero evidence. The result is a phage-level score vector in [0,1]K[0,1]K, where KK is the number of (host, function) 
% combinations with sufficient data coverage.

% A central empirical observation, discovered during data exploration, is that for the gram-type classification task studied
% in Chapter 5, all score dimensions consistently measure gram-positive evidence: phages known to infect gram-positive 
% bacteria achieve mean scores of approximately 0.810.81 across all KK dimensions, while phages infecting gram-negative 
% bacteria achieve mean scores of approximately 0.080.08. This strong and consistent separation is what makes the gram-type 
% task tractable and serves as a proof-of-concept for the full multi-host framework.

